% !TEX program = pdflatex
\documentclass{beamer}

\usetheme{Boadilla}
\usepackage{graphicx}
\usepackage{hyperref}

\title{Early Diffusion Signals for Predicting Reach and Veracity in Social Media Rumor Cascades}
\subtitle{Paper track: thesis presentation}
\author{Jiahao Zhang}
\institute{University of Chicago}
\date{February 2026}

\begin{document}

\section{Background}
\begin{frame}
\titlepage
% NOTES: Open with the core research theme and keep the framing concise. Transition quickly to the agenda.
\end{frame}

\begin{frame}{Outline}
\tableofcontents
% NOTES: Signal the flow: motivation, setup, results, implications, and next steps. Keep timing tight for a 5--10 minute talk.
\end{frame}

\begin{frame}{Motivation: Why Early Warning Matters}
\begin{itemize}
\item \textbf{Timing Problem:} misinformation can scale before reliable verification arrives.
\item \textbf{Example 1 (large reach):} ``did you know that microsoft actually uses dog noses for thumbsticks on the gamecube controller?'' (Twitter16, root \texttt{640644123339431936}).
\item \textbf{Example 2 (smaller reach):} ``president of argentina adopts jewish godson to 'stop him turning into a werewolf' '' (Twitter15, root \texttt{549124560793374720}).
\item \textbf{Early Similarity:} at \(T=60\) min both cascades look very similar in early structure (volume/depth/width around \(65/2/56\sim57\)).
\item \textbf{Outcome Divergence:} final size diverges sharply (\(1453\) vs \(124\)).
\item \textbf{Implication:} this is exactly why early-warning signals are valuable for timely triage.
\end{itemize}
% NOTES: Keep this slide concrete and narrative: same early snapshot, different eventual reach.
\end{frame}

\begin{frame}{Brief Literature}
\begin{itemize}
\item \textbf{Prior Findings:} prior studies find veracity-linked diffusion differences, but mostly from retrospective or full-cascade analysis (e.g., Vosoughi et al., 2018; Del Vicario et al., 2016).
\item \textbf{Open Gap:} it remains unclear whether useful signals appear early enough for actionable intervention.
\item \textbf{This Approach:} we use early observation windows with structural and temporal features beyond early volume alone.
\item \textbf{Scope Today:} this presentation centers the FibVID main analysis and uses Twitter15/16 as a benchmark comparison.
\end{itemize}
% NOTES: Keep the citation visible on the first bullet and keep each bullet title word bolded.
\end{frame}

\section{Data \& Methods}
\begin{frame}{Data}
\begin{itemize}
\item \textbf{Main dataset:} FibVID claim-level propagation records.
\item \textbf{Comparative benchmark:} Twitter15/16 rumor cascade benchmark.
\item \textbf{Main unit of analysis:} one claim-level cascade under a harmonized cascade-compatible representation.
\item \textbf{FibVID label setup:} binary false vs true in the main specification.
\item \textbf{Structure means:} for each cascade or claim-level unit we observe propagation links and timestamps:
\begin{itemize}
\item nodes = tweets/posts in the cascade
\item edges = repost/reply propagation links
\item timestamps = posting time of each node
\end{itemize}
\item \textbf{Why this matters:} harmonized or native propagation structure + timing lets us cut the first \(T\) minutes (or first \(k\) nodes) without peeking into the future.
\end{itemize}
% NOTES: Keep this slide concrete and reproducible: N, labels, and what "full structure" technically contains.
\end{frame}

\begin{frame}{Methods}
\begin{itemize}
\item \textbf{Step 1: Early-window setup} For each cascade, observe only an early prefix: primary size-based windows (first \(k\) posts, holding early volume fixed) and time-based \(T\) as a robustness check.
\item \textbf{Step 2: Feature construction} Compute early features from the partial tree: volume baseline, structural shape, temporal/dynamic features, and structure--tempo interactions.
\item \textbf{Step 3: Task-specific modeling} Same inputs, two targets: reach prediction (regress \(\log(1+\text{final size})\)) and veracity prediction (\textbf{false vs not-false} classification).
\item \textbf{Step 4: Window-by-window evaluation} Repeat across windows to test when early signals become informative and whether patterns differ by task.
\end{itemize}
% NOTES: Walk through pipeline steps once and point to the schematic. Avoid result interpretation on this slide.
\end{frame}

\section{Results}
\begin{frame}{Results (Twitter15/16): Veracity Prediction}
\begin{itemize}
\item \textbf{Signal Exists:} early diffusion features provide above-chance veracity signal (AUC \(>\) 0.5).
\item \textbf{Limited Increment:} gains from richer feature bundles are modest across windows.
\end{itemize}
\begin{columns}[T,onlytextwidth]
\column{0.69\textwidth}
\centering
\IfFileExists{figs/F1K_k_veracity_primary.pdf}{\includegraphics[width=0.99\linewidth]{figs/F1K_k_veracity_primary.pdf}}{\small (Missing file: figs/F1K_k_veracity_primary.pdf)}
\column{0.29\textwidth}
\IfFileExists{tables/twitter15_16_veracity_results.tex}{\centering
\small
\begin{tabular}{lcc}
\hline
Model & AUC (k=60) & SE \\
\hline
Logit & 0.640 & 0.014 \\
RF & 0.633 & 0.012 \\
XGBoost & 0.623 & 0.014 \\
\hline
\end{tabular}
}{\small (Missing file: tables/twitter15_16_veracity_results.tex)}
\end{columns}
% NOTES: Highlight the main directional takeaway without introducing new metrics. Keep language strictly tied to displayed outputs.
\end{frame}

\begin{frame}{Results (Twitter15/16): Reach Prediction}
\begin{itemize}
\item \textbf{Stronger Predictability:} early diffusion is more informative for reach than for veracity.
\item \textbf{Window Effect:} improvements from structure/interaction features become clearer as the window expands.
\end{itemize}
\begin{columns}[T,onlytextwidth]
\column{0.69\textwidth}
\centering
\IfFileExists{figs/F2K_k_reach_primary.pdf}{\includegraphics[width=0.99\linewidth]{figs/F2K_k_reach_primary.pdf}}{\small (Missing file: figs/F2K_k_reach_primary.pdf)}
\column{0.29\textwidth}
\IfFileExists{tables/twitter15_16_reach_results.tex}{\centering
\small
\begin{tabular}{lcc}
\hline
Model & R\textsuperscript{2} (k=60) & SE \\
\hline
OLS & 0.254 & 0.054 \\
RFReg & 0.370 & 0.026 \\
XGBReg & 0.345 & 0.021 \\
\hline
\end{tabular}
}{\small (Missing file: tables/twitter15_16_reach_results.tex)}
\end{columns}
% NOTES: Keep comparison conceptual: same framework, different task. Avoid adding claims beyond the shown Twitter15/16 outputs.
\end{frame}

\begin{frame}{Results: Permutation Test}
\begin{itemize}
\item \textbf{Permutation Check:} we ran a permutation test at \(k=60\) to verify the learned signal is not due to random label assignment.
\end{itemize}
\begin{columns}[T,onlytextwidth]
\column{0.49\textwidth}
\centering
\small Veracity panel
\IfFileExists{figs/F5_permutation_test_signal.pdf}{\includegraphics[width=\linewidth,trim=0 0 428 0,clip]{figs/F5_permutation_test_signal.pdf}}{\small (Missing file: figs/F5_permutation_test_signal.pdf)}
\column{0.49\textwidth}
\centering
\small Reach panel
\IfFileExists{figs/F5_permutation_test_signal.pdf}{\includegraphics[width=\linewidth,trim=428 0 0 0,clip]{figs/F5_permutation_test_signal.pdf}}{\small (Missing file: figs/F5_permutation_test_signal.pdf)}
\end{columns}
\end{frame}

\section{Discussion}
\begin{frame}{Discussion}
\begin{itemize}
\item \textbf{Practical Use:} supports earlier triage, monitoring priority, and resource allocation workflows.
\item \textbf{Conceptual Signal:} early cascade structure appears informative beyond simple early volume counts.
\item \textbf{Limitations:} dataset and window-design choices may affect transferability and operational fit.
\item \textbf{Risks to Monitor:} leakage pathways, heavy-tailed outcomes, and sensitivity to evaluation setup.
\end{itemize}
% NOTES: Balance practical implications with caution on limits. Frame this as evidence for utility, not a final deployment claim.
\end{frame}

\section{Next Steps}
\begin{frame}{Future Work}
\begin{itemize}
\item \textbf{New Primary Dataset:} move to FibVID as the main empirical dataset for the next stage, and keep Twitter15/16 as a robustness check because it is older.
\item \textbf{Reproducible Replication:} rebuild the same early-window pipeline on FibVID (same window definitions, feature groups, and metrics) so results are directly comparable.
\item \textbf{Focused Validation:} feature ablation (stepwise remove/add feature groups to quantify each bundle's incremental contribution).
\end{itemize}
% NOTES: Mark all items as planned work and keep non-numeric. Position this as the roadmap after current Twitter15/16 findings.
\end{frame}

\section{Q\&A}
\begin{frame}{Q\&A + GitHub Repo}
\begin{itemize}
\item Code repository: \href{https://github.com/JiahaoZhang2001/Early-diffusion-signals}{github.com/JiahaoZhang2001/Early-diffusion-signals}
\item Questions?
\end{itemize}
% NOTES: Close with the repository link and invite discussion. Keep this slide minimal and readable.
\end{frame}

\end{document}
